\documentclass[11pt]{article}
\usepackage[margin=25mm]{geometry}
\usepackage{amsmath,amssymb,amsthm}
\usepackage[T1]{fontenc}
\usepackage[colorlinks=true,urlcolor=blue]{hyperref}
\newtheorem{theorem}{Theorem}
\setlength{\parindent}{0pt}\setlength{\parskip}{5pt}
\setlength{\emergencystretch}{3em}
\begin{document}
\begin{center}{\Large The leading coefficient obstruction to an Eisenstein generator}\\[5pt]{\large Conjecture 00000003898}\\[5pt]zbw123456 \quad 14 October 2026\end{center}

Fix a prime $p$. Write $\mathbb Z_p$ for the $p$-adic integers and $v_p(p)=1$ for the normalized valuation. A polynomial $E\in\mathbb Z_p[q]$ is Eisenstein at $(p)$ only if its leading coefficient does not belong to $(p)$; all lower coefficients belong to $(p)$ and the constant coefficient does not belong to $(p)^2$.
\begin{theorem}
There is no polynomial $E\in\mathbb Z_p[q]$ that is Eisenstein at $(p)$ and whose leading coefficient has valuation one.
\end{theorem}
\begin{proof}
Let $a$ be its leading coefficient. The valuation condition $v_p(a)=1$ implies $a\ne0$ and $a\in(p)$. Eisenstein's leading-coefficient condition says $a\notin(p)$. These contradict each other.
\end{proof}
This disproves the existence assertion about the generator before any property of the map $\theta$ is needed. Requiring the polynomial also to generate a kernel or characterize an admissible parameter cannot make a nonexistent polynomial exist. In the usual definition it is the \emph{constant} coefficient of an Eisenstein polynomial that has valuation exactly one; the leading coefficient is a unit. Changing ``leading'' to ``constant'' would be a different problem. Both official language versions say leading coefficient.

\textbf{Lean correspondence.} The project uses Mathlib's $p$-adic integers, Eisenstein predicate, ideal span and valuation. The library theorem \texttt{mem\_span\_pow\_iff\_le\_valuation} converts the valuation condition into ideal membership. \texttt{leading\_valuation\_ne\_one} proves the contradiction for every prime $p$ and polynomial; \texttt{conjecture\_false} negates the required existential statement. The type \texttt{Fact p.Prime} expresses the prime hypothesis and is not a custom axiom. No prism construction or conjectured theorem is assumed.


\textbf{Verification.} Lean 4.19.0 with Mathlib pinned to
\texttt{c44e0c8ee63ca166450922a373c7409c5d26b00b}.
The accompanying project and verification transcript give the exact theorem and axiom audit. No proof placeholders, native evaluator proof bypass, or custom axiom declarations are used. This is a prepared submission; official acceptance has not been claimed.
\end{document}
